\subsection{ARCH-136-FREE-MAP-HOMOTOPY: prescribed parallel maps from the free source}
\label{sec:arch136freemaphomotopy}

\paragraph{Statement.}
Let $k$ be any commutative ring and $C\in\mathcal A_k$.  Choose
$\mu_C(q)$ such that $C_*^{w,q}$ is horizontally acyclic for
$w\geq\mu_C(q)$.  Let $H$ be a simplicial strict differential PD
$S_k$-algebra obtained from the constant $S_k$ by successive free
horizontal attachments
\[
 \mathbb S^r\otimes E_q\ \longrightarrow\
 \mathbb D^{r+1}\otimes E_q,
 \qquad E_q=(k a\xrightarrow{d}k b),\quad |a|=q,
 \quad |b|=q+1,\quad da=b,                                \tag{136.FM1}
\]
where the common weight of each pair is $w\geq\mu_C(q)$.
The arrows in (136.FM1) mean the corresponding free differential
PD algebra maps, used in pushouts under $S_k$.  Arbitrary sets of
cells can be attached at one stage, and filtered unions of stages
are allowed.  In particular the pure source of
ARCH-136-DISK-SUPPORT has this property for every one of its given
targets.

Any two strict $S_k$-preserving maps $f_0,f_1:H\to C$ admit a single
strict coefficient-preserving simplicial homotopy
\[
                  h:H\longrightarrow P_{\Delta[1]}C,
                  \qquad e_i h=f_i.                       \tag{136.FM2}
\]
No admissibility of $H$ is assumed.  Since $P_{\Delta[1]}C$ is
admissible, (136.FM2) is also a simultaneous admissible factorization
of these two prescribed maps under $H$.

More generally, any finite nonempty family of strict maps from $H$
to the same $C$ factors through an admissible finite-path cotensor
$P_JC$, using its successive vertex evaluations.  Therefore the
higher-completion obstruction of AQ-136-HIGHER-COMPLETION cannot
occur for prescribed parallel maps from a source satisfying
(136.FM1) and its weight bounds.

\paragraph{Proof.}
Fix one attached cell's weight $w$ and form degree $q$, and put
$F=C_*^{w,q}$.  By hypothesis $NF$ is acyclic, including $H_0$.
The endpoint restriction
\[
        e:F^{\Delta[1]}\longrightarrow F\oplus F
\]
is degreewise surjective: normalized relative chains of
$\Delta[n]\times\{0,1\}\subset\Delta[n]\times\Delta[1]$
are free, and a prescribed chain map extends one cell at a time
into the acyclic complex $NF$.  This is the endpoint argument
in AQ-136-COTENSOR and ARCH-136-PATHS.  Both its source and target
are horizontally acyclic; for the source, contract the interval
in the cotensor.  Consequently its kernel $K$ has acyclic
normalization, by exactness of normalization and the long exact
homology sequence.  No flatness hypothesis on $k$ is used.

Suppose a homotopy has already been defined on the stage before
the chosen cell.  The boundary of that cell has a prescribed image
in $P_{\Delta[1]}C$.  Its two endpoints on the entire cell are
prescribed by $f_0,f_1$.  For its lower vertical generator, this is
a square of simplicial $k$-modules
\[
 \begin{array}{ccc}
 \mathbb S^r&\longrightarrow&F^{\Delta[1]}\\
 \big\downarrow&&\big\downarrow e\\
 \mathbb D^{r+1}&\longrightarrow&F\oplus F.
 \end{array}                                             \tag{136.FM3}
\]
It has a lift.  Explicitly, on normalized chains choose a lift $y$
of the prescribed image of the degree-$(r+1)$ generator along the
surjection $N_{r+1}e$.  If $z$ is the already prescribed image of
its boundary, then $\partial y-z\in N_rK$ is a cycle.
Acyclicity supplies $v\in N_{r+1}K$ with
$\partial v=\partial y-z$.  Replacing $y$ by $y-v$ gives the
required lift.  This also works when $r=0$, precisely because
$H_0(NK)=0$.  Dold--Kan converts it back to the lift in (136.FM3).

Send the partner generator of form degree $q+1$ to the vertical
derivative of this lift.  Its boundary and endpoints then agree
automatically with the prescribed data because every existing
map commutes with $d$.  The free differential PD universal property
extends these two assignments across the algebra attachment.
Thus a homotopy on a stage extends across the next stage, with its
endpoints fixed.  Cells attached simultaneously can be treated
independently over their common previous stage.

Begin on $S_k$ with constant coefficient paths.  Perform the lifting
successively at every stage and take the compatible union at each
limit stage.  The colimit universal property gives a strict map
to the fixed target $P_{\Delta[1]}C$; no commutation of a cotensor
with a colimit is being asserted.  All endpoints and coefficient
identities are retained.  This proves (136.FM2).

For a finite list $f_1,\ldots,f_s$, choose one homotopy between each
successive pair.  Glue its $s-1$ intervals at their matching
endpoints.  The resulting sections give $H\to P_JC$ with the
specified vertex evaluations.  When $s=1$, use a constant path on
one interval.  ARCH-136-PATHS supplies admissibility of $P_JC$.
This proves the final assertion.\hfill$\square$

\paragraph{What the secondary-period test now says.}
For two maps into the standard theory $B$, the proposed mod-$p$
secondary period on a simplicial sphere cannot distinguish these
maps from a source of the form (136.FM1).  They already factor
through the single admissible theory $P_{\Delta[1]}B$.
Equivalently, the necessary derived-coefficient comparison of
AQ-136-DERIVED-COMPATIBILITY holds for that pair, in its uniform
weight range and for every modulus.  This conclusion uses the
whole cellwise lifting proof, not an unproved assertion that the
individual divided-power expression has an integral primitive.

The result concerns a common target.  It does not supply a common
admissible factorization for maps $H\to C$ and $H\to D$ when
$C,D$ are different objects of unknown connected-component
relation.  Nor does it make an inverse refinement preserving low
free generators possible: ARCH-136-LOW-REFINEMENT still applies.
The remaining forward construction must accommodate different
target theories, or first place them in one component.  Thus the
specific prescribed parallel-map obstruction has been removed for
the actual free source, while integral connectedness remains open.
